\chapter{La suma y la resta}\label{ch:g3:addsub}

Sumar es juntar; restar es quitar o
averiguar lo que falta. Este capítulo entrena los dos métodos en columna
--- con sus famosas llevadas y préstamos --- y los trucos de cálculo mental
que a menudo hacen innecesarias las columnas.

\section{La suma en columna}

\begin{method}[Sumar en columna]\label{met:g3:addsub:add}
\begin{enumerate}
  \item escribe los números uno debajo del otro, unidades bajo unidades y
        decenas bajo decenas (alinéalos por la \emph{derecha});
  \item \omterm{def:g1:addition:def}{suma} las unidades; si la \omterm{def:g1:addition:def}{suma} pasa de $9$, escribe su cifra de las unidades
        y \emph{lleva} la decena a la columna siguiente (un $1$ pequeño);
  \item sigue con las decenas, las centenas, \dots, sumando cada vez la
        llevada.
\end{enumerate}
\end{method}

\begin{example}\label{ex:g3:addsub:add}
Calcula $457 + 368$:
\[
\begin{array}{r}
{\scriptstyle 1}\ \ {\scriptstyle 1}\ \ \phantom{0} \\[-3pt]
4\,5\,7 \\
+\ 3\,6\,8 \\
\hline
8\,2\,5
\end{array}
\]
Unidades: $7 + 8 = 15$: escribe $5$ y lleva $1$. Decenas:
$5 + 6 + 1 = 12$: escribe $2$ y lleva $1$. Centenas: $4 + 3 + 1 = 8$.
Resultado: $825$.
\end{example}

\begin{example}[Sumar de cabeza]\label{ex:g3:addsub:mental}
Busca números amigos antes de recurrir a las columnas:
\begin{itemize}
  \item \emph{formar una decena}: $47 + 8 = 47 + 3 + 5 = 50 + 5 = 55$;
  \item \emph{saltar por partes}: $256 + 130 = 256 + 100 + 30 = 386$;
  \item \emph{reordenar}: $17 + 59 + 3 = (17 + 3) + 59 = 79$.
\end{itemize}
\end{example}

\begin{omfigure}
\begin{tikzpicture}[scale=1.05]
  \draw[-Stealth] (-0.2,0) -- (10.4,0);
  \foreach \x/\l in {0.56/{256}, 5.56/{356}, 7.06/{386}}
    { \draw (\x,-0.08) -- (\x,0.08);
      \node[below, font=\small] at (\x,-0.1) {$\l$}; }
  \draw[-Stealth, omDef, thick] (0.56,0.25) .. controls (2.8,1.15) ..
    (5.56,0.25);
  \node[omDef, font=\small] at (3.05,1.05) {$+100$};
  \draw[-Stealth, omThm, thick] (5.56,0.25) .. controls (6.3,0.8) ..
    (7.06,0.25);
  \node[omThm, font=\small] at (6.3,0.85) {$+30$};
\end{tikzpicture}

{\small La \omterm{def:g1:addition:def}{suma} mental como saltos en la recta numérica: $256 + 130$ en
dos saltos, primero las centenas y luego las decenas.}
\end{omfigure}

\section{La resta en columna}

\begin{method}[Restar en columna]\label{met:g3:addsub:sub}
\begin{enumerate}
  \item escribe arriba el número mayor, alineado por la derecha;
  \item resta columna a columna, empezando por las unidades;
  \item si la cifra de arriba es demasiado pequeña, pide \emph{prestado}: toma una del
        lugar siguiente de la fila de arriba (vale diez en la columna
        actual);
  \item \emph{comprueba}: el resultado más el número restado debe
        devolverte el número de arriba.
\end{enumerate}
\end{method}

\begin{example}\label{ex:g3:addsub:sub}
Calcula $603 - 147$:
\[
\begin{array}{r}
6\,0\,3 \\
-\ 1\,4\,7 \\
\hline
4\,5\,6
\end{array}
\]
Unidades: $3 - 7$ no se puede; pide prestado en cadena: $603$ pasa a ser
$5$ centenas, $9$ decenas y $13$ unidades. Entonces $13 - 7 = 6$;
decenas, $9 - 4 = 5$; centenas, $5 - 1 = 4$. Resultado: $456$.
Comprobación: $456 + 147 = 603$.
\end{example}

\begin{omfigure}
\begin{tikzpicture}[scale=0.82]
  % antes: 6 centenas, 0 decenas, 3 unidades
  \node[font=\small] at (2.5,2.25) {$603$ antes del préstamo};
  \foreach \i in {0,...,5}
    \draw[omProp, fill=omProp!25] (0.62*\i,0.7) rectangle
      (0.62*\i+0.44,1.7);
  \node[below, font=\footnotesize] at (1.7,0.55) {$6$ centenas};
  \node[font=\footnotesize] at (4.35,1.2) {$0$ decenas};
  \foreach \i in {0,1,2} \fill[omDef] (5.4+0.4*\i,1.2) circle (0.13);
  \node[below, font=\footnotesize] at (5.8,0.55) {$3$ unidades};
  % despues: 5 centenas, 9 decenas, 13 unidades
  \begin{scope}[yshift=-3.3cm]
  \node[font=\small] at (2.5,2.25) {$603$ después del préstamo};
  \foreach \i in {0,...,4}
    \draw[omProp, fill=omProp!25] (0.62*\i,0.7) rectangle
      (0.62*\i+0.44,1.7);
  \node[below, font=\footnotesize] at (1.4,0.55) {$5$ centenas};
  \foreach \i in {0,...,8}
    \draw[omThm, very thick] (3.5+0.22*\i,0.7) -- (3.5+0.22*\i,1.7);
  \node[below, font=\footnotesize] at (4.4,0.55) {$9$ decenas};
  \foreach \i in {0,...,6} \fill[omDef] (6.1+0.4*\i,1.45) circle (0.13);
  \foreach \i in {0,...,5} \fill[omDef] (6.1+0.4*\i,0.95) circle (0.13);
  \node[below, font=\footnotesize] at (7.3,0.55) {$13$ unidades};
  \end{scope}
\end{tikzpicture}

{\small El préstamo de $603 - 147$ en dibujos: una centena se desata
en diez decenas y una de esas decenas en diez unidades. La misma cantidad
con ropa nueva --- $6\,\vert\,0\,\vert\,3$ pasa a ser
$5\,\vert\,9\,\vert\,13$, y ya todas las columnas se pueden restar.}
\end{omfigure}

\begin{example}[La resta como la parte que falta]\label{ex:g3:addsub:missing}
«Lila tiene $35$ canicas y Tomás tiene $52$: ¿cuántas tiene Tomás de más?» es
la resta $52 - 35$. De cabeza, cuenta hacia \emph{delante} desde $35$:
del $35$ al $40$ hay $5$ y del $40$ al $52$ hay $12$; en total,
$5 + 12 = 17$. Contar hacia delante suele ser más fácil que pedir prestado.
\end{example}

\begin{omfigure}
\begin{tikzpicture}[scale=1.05]
  \draw[-Stealth] (-0.2,0) -- (10.4,0);
  \foreach \x/\l in {1.0/{35}, 3.5/{40}, 9.5/{52}}
    { \draw (\x,-0.08) -- (\x,0.08);
      \node[below, font=\small] at (\x,-0.1) {$\l$}; }
  \draw[-Stealth, omDef, thick] (1.0,0.25) .. controls (2.2,0.9) ..
    (3.5,0.25);
  \node[omDef, font=\small] at (2.25,0.92) {$+5$};
  \draw[-Stealth, omThm, thick] (3.5,0.25) .. controls (6.5,1.2) ..
    (9.5,0.25);
  \node[omThm, font=\small] at (6.5,1.1) {$+12$};
\end{tikzpicture}

{\small El cálculo de $52 - 35$ contando hacia delante: primero hasta el amable
$40$ y después hasta $52$. La respuesta es el total de los saltos:
$5 + 12 = 17$.}
\end{omfigure}

\section{Órdenes de magnitud}

\begin{method}[Estimar antes de calcular]\label{met:g3:addsub:estimate}
Antes de cualquier cálculo en columna, redondea los números a otros amables y
calcula una \emph{estimación}. Después de calcular, compara: un resultado lejos
de la estimación avisa de un error.
\end{method}

\begin{example}
Para $487 + 312$, la estimación es $500 + 300 = 800$. El resultado en
columna, $799$, está cerca: es verosímil. Si hubieras obtenido $1\,099$
(por una mala alineación), la estimación te habría avisado.
\end{example}

\section{Ejercicios}

\begin{exercise}[$\star$]\label{exo:g3:addsub:1}
Calcula en columna: $346 + 253$; \quad $478 + 265$; \quad
$1\,536 + 2\,876$.
\end{exercise}

\begin{exercise}[$\star$]\label{exo:g3:addsub:2}
Calcula en columna y comprueba después con una \omterm{def:g1:addition:def}{suma}: $587 - 234$; \quad
$805 - 348$; \quad $1\,000 - 463$.
\end{exercise}

\begin{exercise}[$\star$]\label{exo:g3:addsub:3}
Calcula de cabeza formando decenas: $38 + 7$; \quad $56 + 9$; \quad
$45 + 28$ (\omterm{def:g1:addition:def}{suma} $30$ y quita $2$).
\end{exercise}

\begin{exercise}[$\star$]\label{exo:g3:addsub:4}
Calcula de cabeza por saltos: $340 + 220$; \quad $675 + 210$; \quad
$512 - 300$; \quad $840 - 220$.
\end{exercise}

\begin{exercise}[$\star$]\label{exo:g3:addsub:5}
Calcula $63 - 47$ contando hacia delante (como en el
\cref{ex:g3:addsub:missing}) y comprueba en columna.
\end{exercise}

\begin{exercise}[$\star$]\label{exo:g3:addsub:6}
Estima primero (redondeando a centenas) y después calcula con exactitud:
$492 + 218$; \quad $913 - 388$.
\end{exercise}

\begin{exercise}[$\star$]\label{exo:g3:addsub:7}
Un colegio tiene $348$ niñas y $296$ niños. ¿Cuántos alumnos hay en total?
Escribe la operación, calcula y responde con una frase.
\end{exercise}

\begin{exercise}[$\star$]\label{exo:g3:addsub:8}
Un libro tiene $224$ páginas y Nora ha leído $178$. ¿Cuántas páginas le
quedan por leer?
\end{exercise}

\begin{exercise}[$\star$]\label{exo:g3:addsub:9}
Completa los huecos (trabaja columna a columna):
\[
\begin{array}{r}
3\,?\,7 \\
+\ ?\,2\,? \\
\hline
8\,6\,2
\end{array}
\]
¿Hay una sola manera de rellenar los tres huecos? Explícalo.
\end{exercise}

\begin{exercise}[$\star\star$]\label{exo:g3:addsub:10}
Una panadería vendió $145$ cruasanes por la mañana y $89$ por la
tarde; al cerrar quedaron $17$ cruasanes sin vender. ¿Cuántos
cruasanes se hornearon ese día? (Dos pasos.)
\end{exercise}

\begin{exercise}[$\star\star$]\label{exo:g3:addsub:11}
Teo dice: «para calcular $500 - 199$, calculo $500 - 200$ y después
sumo $1$». ¿Tiene razón? Calcula de las dos maneras y explica su truco.
\end{exercise}
